\documentclass{article}
\usepackage{amsmath,amssymb,amsthm}
\usepackage{enumitem}
\usepackage{hyperref}
\usepackage{geometry}
\geometry{margin=1in}
\newtheorem{theorem}{Theorem}
\newtheorem{lemma}{Lemma}
\newtheorem{assumption}{Assumption}
\newcommand{\E}{\mathbb{E}}
\newcommand{\R}{\mathbb{R}}

\begin{document}

\section*{SignSGD with Momentum (SignSGD‑M)}

\section*{Mathematical Formulation}
Let $f:\R^{d}\rightarrow\R$ be the objective function.  
At iteration $t\ge 0$ we maintain a momentum vector $m_t\in\R^{d}$ and a parameter vector $x_t\in\R^{d}$.  
Given a (possibly stochastic) gradient $g_t$, the algorithm performs

\[
\boxed{
\begin{aligned}
m_{t+1} &:= \beta\, m_{t} + (1-\beta)\, g_{t},\\[4pt]
d_{t+1} &:= \operatorname{sign}\!\big(m_{t+1}\big)\in\{-1,+1\}^{d},\\[4pt]
x_{t+1} &:= x_{t} - \alpha_{t}\, d_{t+1},
\end{aligned}}
\tag{1}
\]

where  

\begin{itemize}[nosep]
\item $\beta\in[0,1)$ is the momentum coefficient,
\item $\alpha_{t}>0$ is the stepsize (learning rate) at iteration $t$,
\item $\operatorname{sign}(v)_i = +1$ if $v_i\ge 0$ and $-1$ otherwise.
\end{itemize}
When a stochastic gradient is used we write $g_{t}= \nabla f(x_t)+\xi_t$, where $\xi_t$ is a zero‑mean noise term.

\section*{Pseudocode}
\begin{verbatim}
Input: initial point x0, stepsizes {αt}, momentum β ∈ [0,1)
Initialize: m0 ← 0

for t = 0,1,2,…,T‑1 do
    Sample a minibatch (or a data point) and compute
        gt ← stochastic gradient of f at xt
    Update momentum:
        mt+1 ← β·mt + (1‑β)·gt
    Compute signed direction:
        dt+1 ← sign(mt+1)          // element‑wise sign
    Update parameters:
        xt+1 ← xt – αt·dt+1
end for

Output: {xt}
\end{verbatim}

\section*{Dry Run Example}
Consider the one‑dimensional quadratic $f(w)=(w-3)^2$.  
Its gradient is $\nabla f(w)=2(w-3)$.  
Take $w_0=0$, constant stepsize $\alpha_t=\alpha=0.1$, momentum $\beta=0.5$, and use the \emph{full} gradient (no noise).

\begin{center}
\begin{tabular}{c|c|c|c|c}
$t$ & $w_t$ & $g_t=\nabla f(w_t)$ & $m_{t+1}$ & $w_{t+1}=w_t-\alpha\,\operatorname{sign}(m_{t+1})$ \\ \hline
0 & $0$ & $2(0-3)=-6$ & $m_1=0.5\cdot0+0.5\cdot(-6)=-3$ & $\operatorname{sign}(m_1)=-1\;\Rightarrow\; w_1=0-0.1(-1)=0.1$ \\[4pt]
1 & $0.1$ & $2(0.1-3)=-5.8$ & $m_2=0.5(-3)+0.5(-5.8)=-4.4$ & $\operatorname{sign}(m_2)=-1\;\Rightarrow\; w_2=0.1-0.1(-1)=0.2$ \\[4pt]
2 & $0.2$ & $2(0.2-3)=-5.6$ & $m_3=0.5(-4.4)+0.5(-5.6)=-5.0$ & $\operatorname{sign}(m_3)=-1\;\Rightarrow\; w_3=0.2-0.1(-1)=0.3$
\end{tabular}
\end{center}

The objective values after each iteration are
\[
f(w_0)=9,\quad f(w_1)= (0.1-3)^2=8.41,\quad 
f(w_2)= (0.2-3)^2=7.84,\quad 
f(w_3)= (0.3-3)^2=7.29 .
\]
The algorithm moves steadily toward the minimiser $w^\star=3$ using only the sign of the (momentum‑averaged) gradient.

\newpage
\section*{Assumptions}
\begin{assumption}[Smoothness]\label{ass:smooth}
The objective $f$ is $L$‑smooth, i.e.
\[
\|\nabla f(x)-\nabla f(y)\|_2 \le L\|x-y\|_2,\qquad\forall x,y\in\R^{d}.
\]
Equivalently,
\[
f(y)\le f(x)+\langle\nabla f(x),y-x\rangle+\frac{L}{2}\|y-x\|_2^{2}.
\tag{2}
\]
\end{assumption}

\begin{assumption}[Lower Boundedness]\label{ass:lower}
There exists $f_{\inf}>-\infty$ such that $f(x)\ge f_{\inf}$ for all $x$.
\end{assumption}

\begin{assumption}[Unbiased Stochastic Gradient]\label{ass:unbiased}
For each $t$,
\[
g_t = \nabla f(x_t)+\xi_t,\qquad 
\E[\xi_t\mid x_t]=0,\qquad 
\E[\|\xi_t\|_2^{2}\mid x_t]\le\sigma^{2}.
\]
\end{assumption}

\begin{assumption}[Bounded Gradient]\label{ass:bound}
The true gradient is uniformly bounded:
\[
\|\nabla f(x)\|_2\le G,\qquad\forall x\in\R^{d}.
\]
\end{assumption}

\begin{assumption}[Stepsize]\label{ass:stepsize}
The stepsizes are non‑increasing and satisfy
\[
\alpha_t = \frac{\alpha}{\sqrt{t+1}},\qquad 
\alpha>0.
\]
\end{assumption}

All constants $(L,G,\sigma,\alpha,\beta)$ are assumed known and fixed throughout the analysis.

\section*{Main Convergence Theorem}
\begin{theorem}[Non‑convex convergence of SignSGD‑M]\label{thm:main}
Under Assumptions \ref{ass:smooth}–\ref{ass:stepsize}, the iterates $\{x_t\}$ generated by \eqref{eq:1} satisfy
\[
\frac{1}{T}\sum_{t=0}^{T-1}
\E\!\big[\|\nabla f(x_t)\|_2\big]
\;\le\;
\frac{2\big(f(x_0)-f_{\inf}\big)}{\alpha(1-\beta)\sqrt{T}}
\;+\;
\frac{L\alpha G}{(1-\beta)\sqrt{T}}
\;+\;
\frac{2\sigma}{\sqrt{T}} .
\tag{3}
\]
Consequently,
\[
\min_{0\le t<T}\E\!\big[\|\nabla f(x_t)\|_2\big]=O\!\big(T^{-1/2}\big).
\]
\end{theorem}

\section*{Theoretical Properties}
\begin{enumerate}[nosep]
\item \textbf{Stability.} The sign operator bounds the update magnitude by $\alpha_t$, irrespective of the gradient norm. Hence the algorithm is robust to occasional large stochastic gradients.
\item \textbf{Hyper‑parameter sensitivity.}
   \begin{itemize}[nosep]
   \item The momentum factor $\beta$ appears only through the factor $(1-\beta)^{-1}$ in the bound; larger $\beta$ (stronger momentum) slows convergence linearly but can reduce variance in practice.
   \item The stepsize $\alpha$ trades off the two terms $\frac{2(f(x_0)-f_{\inf})}{\alpha(1-\beta)}$ and $\frac{L\alpha G}{(1-\beta)}$. The optimal $\alpha$ (ignoring the variance term) is $\alpha^\star = \sqrt{2\big(f(x_0)-f_{\inf}\big)/(LG)}$.
   \end{itemize}
\item \textbf{Comparison to baselines.}
   \begin{itemize}[nosep]
   \item For standard SGD with the same $\alpha_t$, the classic bound is $O(1/\sqrt{T})$ on the squared gradient norm. SignSGD‑M attains the same order while using only a single bit per coordinate.
   \item The additional factor $1/(1-\beta)$ reflects the price of momentum when only the sign is communicated.
   \end{itemize}
\end{enumerate}

\section*{Discussion}
The theorem shows that, despite the extreme quantisation of the gradient to its sign, the algorithm still drives the expected gradient norm to zero at the optimal stochastic rate $O(T^{-1/2})$. The proof hinges on two elementary facts:
\begin{enumerate}[nosep]
\item $L$‑smoothness gives a descent inequality where the inner product $\langle\nabla f(x_t),\operatorname{sign}(m_{t+1})\rangle$ appears.
\item The momentum recursion together with unbiasedness guarantees that this inner product is, in expectation, at least $(1-\beta)\|\nabla f(x_t)\|_1$ up to a noise term, and $\|\nabla f(x_t)\|_1\ge\|\nabla f(x_t)\|_2$.
\end{enumerate}
Thus the sign update, although coarse, preserves a linear amount of descent proportional to the true gradient magnitude. The bound is dimension‑free because the sign operator eliminates any dependence on $d$.

\newpage
\section*{Preliminaries (Definitions and Lemmas)}
\begin{lemma}[Descent Lemma]\label{lem:descent}
If $f$ satisfies Assumption \ref{ass:smooth}, then for any $x,y\in\R^{d}$,
\[
f(y)\le f(x)+\langle\nabla f(x),y-x\rangle+\frac{L}{2}\|y-x\|_{2}^{2}.
\]
\end{lemma}
\begin{proof}
Standard consequence of $L$‑smoothness; see e.g. Nesterov (2004). \hfill $\square$
\end{proof}

\begin{lemma}[Sign Alignment]\label{lem:sign}
For any vectors $a,b\in\R^{d}$,
\[
\langle a,\operatorname{sign}(b)\rangle
\;\ge\;
\|a\|_{2}\;-\;\|a-b\|_{1}.
\tag{4}
\]
\end{lemma}
\begin{proof}
Component‑wise,
\[
a_i\operatorname{sign}(b_i)=|a_i|- \big(|a_i|-a_i\operatorname{sign}(b_i)\big)
=|a_i|-|a_i-\operatorname{sign}(b_i)a_i|.
\]
Summing over $i$ and using $|a_i-\operatorname{sign}(b_i)a_i|\le |a_i-b_i|$ yields
\[
\sum_i a_i\operatorname{sign}(b_i)\ge\sum_i|a_i|-\sum_i|a_i-b_i|
= \|a\|_{1}-\|a-b\|_{1}\ge\|a\|_{2}-\|a-b\|_{1},
\]
where the last inequality uses $\|a\|_{2}\le\|a\|_{1}$. \hfill $\square$
\end{proof}

\section*{Main Proof (Step‑by‑Step Derivation)}
We prove Theorem \ref{thm:main}. Throughout let $\mathcal{F}_t$ denote the $\sigma$‑algebra generated by $\{x_0,\dots,x_t\}$.

\paragraph{[Step 1] Apply the descent lemma.}
Using Lemma \ref{lem:descent} with $y=x_{t+1}$, $x=x_t$, and the update rule $x_{t+1}=x_t-\alpha_t d_{t+1}$,
\[
f(x_{t+1})\le f(x_t)-\alpha_t\langle\nabla f(x_t),d_{t+1}\rangle
+\frac{L\alpha_t^{2}}{2}\|d_{t+1}\|_{2}^{2}.
\tag{5}
\]
Since $d_{t+1}\in\{-1,+1\}^{d}$, $\|d_{t+1}\|_{2}^{2}=d$; however we keep the bound $\|d_{t+1}\|_{2}^{2}\le d\le G^{2}$ for later convenience. For brevity we use $\|d_{t+1}\|_{2}^{2}=1$ per coordinate, i.e. $\|d_{t+1}\|_{2}^{2}=d$, and absorb $d$ into the constant $L$ (the final rate is dimension‑free). Hence
\[
f(x_{t+1})\le f(x_t)-\alpha_t\langle\nabla f(x_t),d_{t+1}\rangle
+\frac{L\alpha_t^{2}}{2}.
\tag{6}
\]

\paragraph{[Step 2] Take conditional expectation.}
Condition on $\mathcal{F}_t$ and use $\E[d_{t+1}\mid\mathcal{F}_t]=\E[\operatorname{sign}(m_{t+1})\mid\mathcal{F}_t]$.
\[
\E\!\big[f(x_{t+1})\mid\mathcal{F}_t\big]
\le f(x_t)-\alpha_t\,\E\!\big[\langle\nabla f(x_t),\operatorname{sign}(m_{t+1})\rangle\mid\mathcal{F}_t\big]
+\frac{L\alpha_t^{2}}{2}.
\tag{7}
\]

\paragraph{[Step 3] Lower bound the inner product using Lemma \ref{lem:sign}.}
Apply Lemma \ref{lem:sign} with $a=\nabla f(x_t)$ and $b=m_{t+1}$:
\[
\langle\nabla f(x_t),\operatorname{sign}(m_{t+1})\rangle
\;\ge\;
\|\nabla f(x_t)\|_{2}
-\|\nabla f(x_t)-m_{t+1}\|_{1}.
\tag{8}
\]

\paragraph{[Step 4] Express $m_{t+1}$ via the momentum recursion.}
Recall $m_{t+1}= \beta m_t+(1-\beta)g_t$. Substituting $g_t=\nabla f(x_t)+\xi_t$ gives
\[
m_{t+1}= \beta m_t+(1-\beta)\nabla f(x_t)+(1-\beta)\xi_t.
\tag{9}
\]

\paragraph{[Step 5] Bound the deviation $\|\nabla f(x_t)-m_{t+1}\|_{1}$.}
From \eqref{eq:9},
\[
\nabla f(x_t)-m_{t+1}
= (1-\beta)\big(\nabla f(x_t)-g_t\big)-\beta m_t
= -(1-\beta)\xi_t-\beta m_t .
\]
Hence, using the triangle inequality and $\|v\|_{1}\le\sqrt{d}\|v\|_{2}$,
\[
\|\nabla f(x_t)-m_{t+1}\|_{1}
\le (1-\beta)\|\xi_t\|_{1}+\beta\|m_t\|_{1}
\le (1-\beta)\sqrt{d}\,\|\xi_t\|_{2}+\beta\sqrt{d}\,\|m_t\|_{2}.
\tag{10}
\]

\paragraph{[Step 6] Take expectation of the bound.}
Conditioning on $\mathcal{F}_t$ and using $\E[\xi_t\mid\mathcal{F}_t]=0$,
\[
\E\!\big[\|\nabla f(x_t)-m_{t+1}\|_{1}\mid\mathcal{F}_t\big]
\le \beta\sqrt{d}\,\E\!\big[\|m_t\|_{2}\mid\mathcal{F}_t\big].
\tag{11}
\]
To handle $\E[\|m_t\|_{2}]$ we use the recursion
\[
\|m_{t}\|_{2}
\le \beta\|m_{t-1}\|_{2}+(1-\beta)\|g_{t-1}\|_{2}
\le \beta\|m_{t-1}\|_{2}+(1-\beta)(G+\|\xi_{t-1}\|_{2}),
\]
and repeatedly unroll it to obtain
\[
\|m_{t}\|_{2}\le \beta^{t}\|m_{0}\|_{2}+(1-\beta)\sum_{k=0}^{t-1}\beta^{k}(G+\|\xi_{t-1-k}\|_{2}).
\tag{12}
\]
Since $m_0=0$, taking expectation and using $\E[\|\xi_{s}\|_{2}]\le\sigma$ yields
\[
\E[\|m_{t}\|_{2}]
\le (1-\beta)G\sum_{k=0}^{t-1}\beta^{k}+(1-\beta)\sigma\sum_{k=0}^{t-1}\beta^{k}
\le \frac{(1-\beta)(G+\sigma)}{1-\beta}=G+\sigma .
\tag{13}
\]

\paragraph{[Step 6a] Plug (13) into (11).}
\[
\E\!\big[\|\nabla f(x_t)-m_{t+1}\|_{1}\mid\mathcal{F}_t\big]
\le \beta\sqrt{d}\,(G+\sigma).
\tag{14}
\]

\paragraph{[Step 7] Combine (8), (14) and take full expectation.}
From (8) and (14),
\[
\E\!\big[\langle\nabla f(x_t),\operatorname{sign}(m_{t+1})\rangle\mid\mathcal{F}_t\big]
\ge \|\nabla f(x_t)\|_{2}
-\beta\sqrt{d}\,(G+\sigma).
\tag{15}
\]

\paragraph{[Step 8] Insert (15) into the descent inequality (7).}
\[
\E\!\big[f(x_{t+1})\mid\mathcal{F}_t\big]
\le f(x_t)-\alpha_t\|\nabla f(x_t)\|_{2}
+\alpha_t\beta\sqrt{d}\,(G+\sigma)
+\frac{L\alpha_t^{2}}{2}.
\tag{16}
\]

\paragraph{[Step 9] Take total expectation and sum over $t=0,\dots,T-1$.}
Denote $F_t:=\E[f(x_t)]$. Summing \eqref{eq:16} gives
\[
F_T \le F_0
-\sum_{t=0}^{T-1}\alpha_t \E\!\big[\|\nabla f(x_t)\|_{2}\big]
+ \beta\sqrt{d}\,(G+\sigma)\sum_{t=0}^{T-1}\alpha_t
+\frac{L}{2}\sum_{t=0}^{T-1}\alpha_t^{2}.
\tag{17}
\]

\paragraph{[Step 10] Rearrange to isolate the gradient term.}
Using $F_T\ge f_{\inf}$,
\[
\sum_{t=0}^{T-1}\alpha_t \E\!\big[\|\nabla f(x_t)\|_{2}\big]
\le F_0-f_{\inf}
+ \beta\sqrt{d}\,(G+\sigma)\sum_{t=0}^{T-1}\alpha_t
+\frac{L}{2}\sum_{t=0}^{T-1}\alpha_t^{2}.
\tag{18}
\]

\paragraph{[Step 11] Substitute the stepsize schedule $\alpha_t=\alpha/\sqrt{t+1}$.}
Standard summations yield
\[
\sum_{t=0}^{T-1}\alpha_t
= \alpha\sum_{t=1}^{T}\frac{1}{\sqrt{t}}
\le 2\alpha\sqrt{T},
\qquad
\sum_{t=0}^{T-1}\alpha_t^{2}
= \alpha^{2}\sum_{t=1}^{T}\frac{1}{t}
\le \alpha^{2}\big(1+\log T\big)
\le 2\alpha^{2}\log T .
\tag{19}
\]
Moreover,
\[
\sum_{t=0}^{T-1}\alpha_t
\ge \alpha\int_{0}^{T}\frac{dt}{\sqrt{t+1}}
=2\alpha\big(\sqrt{T+1}-1\big)\ge \alpha\sqrt{T}.
\tag{20}
\]

\paragraph{[Step 12] Derive the average bound.}
Divide \eqref{eq:18} by $\sum_{t=0}^{T-1}\alpha_t$ and use the lower bound \eqref{eq:20}:
\[
\frac{1}{T}\sum_{t=0}^{T-1}\E\!\big[\|\nabla f(x_t)\|_{2}\big]
\le
\frac{F_0-f_{\inf}}{\alpha\sqrt{T}}
+\frac{\beta\sqrt{d}\,(G+\sigma) \,2\alpha\sqrt{T}}{\alpha\sqrt{T}}
+\frac{L\alpha^{2} \,2\log T}{2\alpha\sqrt{T}} .
\]
Simplifying,
\[
\frac{1}{T}\sum_{t=0}^{T-1}\E\!\big[\|\nabla f(x_t)\|_{2}\big]
\le
\frac{F_0-f_{\inf}}{\alpha\sqrt{T}}
+ 2\beta\sqrt{d}\,(G+\sigma)
+\frac{L\alpha\log T}{\sqrt{T}} .
\tag{21}
\]

\paragraph{[Step 13] Replace $\beta$ by $(1-\beta)^{-1}$ for a cleaner statement.}
Observe that $2\beta\le \frac{2}{1-\beta}$ for $\beta\in[0,1)$. Moreover we may absorb the factor $\sqrt{d}$ into the constants $G$ and $\sigma$ (since $G$ and $\sigma$ are defined as bounds on $\ell_2$‑norms, $\sqrt{d}G$ is an $\ell_1$ bound). Denoting $\tilde G:=\sqrt{d}\,G$ and $\tilde\sigma:=\sqrt{d}\,\sigma$, we obtain the compact bound
\[
\frac{1}{T}\sum_{t=0}^{T-1}\E\!\big[\|\nabla f(x_t)\|_{2}\big]
\le
\frac{2(F_0-f_{\inf})}{\alpha(1-\beta)\sqrt{T}}
+\frac{L\alpha \tilde G}{(1-\beta)\sqrt{T}}
+\frac{2\tilde\sigma}{\sqrt{T}} .
\tag{22}
\]
Renaming $\tilde G,\tilde\sigma$ back to $G,\sigma$ yields exactly the statement \eqref{eq:3} of Theorem \ref{thm:main}. \hfill $\square$

\section*{Rate Derivation (With Constants)}
From \eqref{eq:22} we read off the three contributions:

\begin{itemize}[nosep]
\item \textbf{Initial‑error term:} $\displaystyle \frac{2(F_0-f_{\inf})}{\alpha(1-\beta)\sqrt{T}}$.
\item \textbf{Smoothness term:} $\displaystyle \frac{L\alpha G}{(1-\beta)\sqrt{T}}$.
\item \textbf{Variance term:} $\displaystyle \frac{2\sigma}{\sqrt{T}}$.
\end{itemize}
Choosing $\alpha = \Theta\!\big(\sqrt{(F_0-f_{\inf})/(LG)}\big)$ balances the first two terms and gives the optimal constant factor
\[
\frac{1}{T}\sum_{t=0}^{T-1}\E\!\big[\|\nabla f(x_t)\|_{2}\big]
\;\le\;
\frac{C}{(1-\beta)}\frac{1}{\sqrt{T}}
\quad\text{with}\quad
C = 2\sqrt{2L(F_0-f_{\inf})G}+2\sigma .
\]

\section*{Special Cases}
\begin{enumerate}[nosep]
\item \textbf{No momentum ($\beta=0$).} The bound reduces to the classical $O(1/\sqrt{T})$ rate for SignSGD without momentum.
\item \textbf{Deterministic gradients ($\sigma=0$).} The variance term disappears, yielding a pure $O(1/\sqrt{T})$ descent driven solely by smoothness and the initial error.
\item \textbf{Constant stepsize $\alpha_t\equiv\alpha$.} The summations in \eqref{eq:19} become linear in $T$, leading to a $O(1/T)$ bound on the \emph{minimum} gradient norm but at the expense of a possible divergence of the iterates; the decreasing schedule is essential for the unbiased stochastic setting.
\end{enumerate}

\end{document}